% !TEX root = ../main.tex
\chapter{Performance Methodology}
\label{ch:perf_methodology}

This part of the report leaves the API and architecture for the
performance-engineering work that shapes \Lucid's
production-readiness. The chapter opens with the methodology
itself: how we measure, what we measure, how we control for
noise, and how we report. Without a disciplined methodology the
performance claims in the following chapters would be unreliable
and the regressions invisible. Establishing the methodology
explicitly --- and revisiting it whenever a measurement looks
suspicious --- is a foundational habit, not a finishing touch.

\section{Why Measurement Discipline}
\label{sec:perf_methodology:why_discipline}

A naive ``run the model, time the iteration'' approach produces
numbers that can vary by 30--50\% between runs on the same
machine. Sources of variance:

\begin{itemize}
  \item \textbf{Thermal throttling}: M3~Max sustains its peak
    only for a few seconds; sustained workloads see clock
    drops of 5--15\%.
  \item \textbf{Background processes}: other apps competing for
    CPU/GPU.
  \item \textbf{Frequency scaling}: the SoC may downclock for
    power management.
  \item \textbf{Cache warming}: the first iteration is slower
    than steady state due to cold caches and JIT compilation.
  \item \textbf{Memory allocator state}: fragmentation history
    affects allocation cost.
\end{itemize}

A measurement that does not control for these is a number
without meaning. The methodology below addresses each.

\section{The Measurement Loop}
\label{sec:perf_methodology:measurement_loop}

The basic \Lucid\ benchmark loop:

\begin{lstlisting}[style=lucid,language=LucidPython]
def bench(fn, *, warmup=10, iters=100):
    # Warm-up: trigger any JIT / cache fill / allocator setup
    for _ in range(warmup):
        fn()
    lucid.runtime.synchronize()

    # Timed iterations
    times = []
    for _ in range(iters):
        lucid.runtime.synchronize()
        t0 = time.perf_counter_ns()
        fn()
        lucid.runtime.synchronize()
        t1 = time.perf_counter_ns()
        times.append((t1 - t0) / 1e6)  # ms

    return {
        "median": statistics.median(times),
        "p05":    sorted(times)[int(iters * 0.05)],
        "p95":    sorted(times)[int(iters * 0.95)],
        "min":    min(times),
        "max":    max(times),
    }
\end{lstlisting}

The key elements:
\begin{enumerate}
  \item \textbf{Warm-up}: drops the first $w$ iterations so
    initial-call overhead does not contaminate the median.
  \item \textbf{Synchronisation around each timed call}: Apple's
    GPU is asynchronous; without sync, \code{perf\_counter\_ns}
    captures only the dispatch time, not the actual execution.
  \item \textbf{Median + percentiles}: median is robust to
    occasional spikes; p05/p95 quantify the spread.
\end{enumerate}

\subsection{Synchronisation primitives}
\label{ssec:perf_methodology:sync}

\code{lucid.runtime.synchronize()} blocks until all submitted
GPU work has completed. Implementation:

\begin{itemize}
  \item On the Metal/MPSGraph path: calls
    \code{[command\_buffer waitUntilCompleted]} on the active
    command buffer and any queued ones.
  \item On the MLX path: calls \code{mx.synchronize()}.
  \item On the CPU/Accelerate path: a no-op (Accelerate is
    synchronous).
\end{itemize}

The function is cheap when there is no outstanding work and
exactly what is needed to make timing well-defined.

\section{Controlling Thermal Variance}
\label{sec:perf_methodology:thermal}

The most insidious source of variance on Apple silicon. Two
strategies:

\subsection{Cooling pause between repeats}
\label{ssec:perf_methodology:cooling}

For multi-configuration sweeps, insert a 30-second pause between
configurations and ignore the first $k$ iterations of each (to
let the chip return to baseline). \Lucid's benchmark harness
implements this automatically when the
\code{LUCID\_BENCH\_THERMAL\_RECOVER=1} env var is set.

\subsection{Sustained-mode benchmarks}
\label{ssec:perf_methodology:sustained}

Some workloads (training a model for hours) actually run at the
thermally-throttled steady state. For these, we run the
benchmark for at least 2 minutes after warm-up and report the
last-minute median, ensuring the thermal state has reached
equilibrium.

\subsection{Power profile control}
\label{ssec:perf_methodology:power_profile}

macOS's power profile (Low Power, Automatic, High Performance)
affects sustained clock targets. For all reported numbers, the
profile is \emph{High Performance} (set via the
``Battery'' / ``Power Adapter'' panes); benchmarks under other
profiles are run only when explicitly relevant (mobile
inference, battery testing).

\section{Profilers Used}
\label{sec:perf_methodology:profilers}

The performance investigations in subsequent chapters draw on
several profilers:

\subsection{Lucid's built-in profiler}
\label{ssec:perf_methodology:lucid_profiler}

The \code{lucid.profiler.Profile} context
(\chref{ch:profiler_determinism}) records per-op timings, memory
allocations, and the MPSGraph compile-vs-execute split.

\begin{lstlisting}[style=lucid,language=LucidPython]
from lucid.profiler import Profile

with Profile(record_shapes=True, with_stack=True) as prof:
    for _ in range(20):
        loss = model(x).sum()
        loss.backward()

print(prof.key_averages().table(sort_by="self_time", row_limit=20))
\end{lstlisting}

The output ranks ops by total time spent; the top entries are
the optimisation targets.

\subsection{Instruments.app}
\label{ssec:perf_methodology:instruments}

Apple's profiler (\code{Instruments.app}, part of Xcode) gives
much finer-grained data: per-Metal-shader dispatch times, GPU
utilisation, memory bandwidth, energy. Used for deep dives when
the \Lucid\ profiler points at an op but the cause is not
obvious from the kernel level.

\subsection{System trace}
\label{ssec:perf_methodology:system_trace}

The \texttt{powermetrics} CLI tool provides 100\,ms-resolution
samples of CPU/GPU clock, power, and temperature. Useful for
spotting thermal throttling in long-running benchmarks.

\section{Reporting Format}
\label{sec:perf_methodology:reporting_format}

Every performance claim in this part of the report follows this
template:

\begin{itemize}
  \item \textbf{Setup}: hardware (M3~Max 16-core CPU, 40-core
    GPU, 64\,GB memory), macOS version (26.0.1), \Lucid\
    version (3.5.0), model + dtype + batch + sequence size.
  \item \textbf{Number}: median of $\geq$\,50 timed iterations
    after $\geq$\,10 warm-up iterations, with explicit p05/p95
    when relevant.
  \item \textbf{Comparison}: if a speedup is claimed, both
    baseline and improved numbers, with both measured in the
    same session to control for thermal drift.
\end{itemize}

Vague claims (``2x faster'') are unacceptable without all four
elements.

\section{Comparison Baselines}
\label{sec:perf_methodology:baselines}

\Lucid\ is compared, where relevant, to two baselines:

\subsection{Reference framework on the same hardware}
\label{ssec:perf_methodology:ref_framework}

For correctness, the \code{lucid/test/\_fixtures/ref\_framework.py}
adapter runs the equivalent computation in a reference ML
framework on the same data. This is what the parity tests use.
For benchmark purposes, the same harness times the reference
framework.

The comparison answers: ``On the same Mac, is \Lucid\ within
20\% of the reference framework's throughput?'' --- the
performance acceptance criterion that has guided every
optimisation.

\subsection{MLX direct}
\label{ssec:perf_methodology:mlx_direct}

For ops that map 1:1 to MLX (most element-wise, reductions,
matmul), we compare \Lucid's call-path overhead vs raw MLX. The
delta is \Lucid's tax for offering autograd and the unified
Tensor abstraction; targets are $<$\,3\,\% overhead on call-
intensive workloads.

\section{Statistical Significance}
\label{sec:perf_methodology:statistical}

For declaring an A/B improvement statistically significant, we
require:

\begin{enumerate}
  \item The medians differ by more than $2\sigma$ of the
    bootstrap distribution of either median.
  \item The improvement is reproducible across 3 separate
    benchmark runs on cold cache (i.e., not just a single
    warm-cache lucky run).
  \item For multi-machine claims, at least 2 different Apple
    silicon variants (M3~Max, M2~Pro, etc.) show the same
    direction of effect.
\end{enumerate}

In practice, most optimisations clear this bar handily (10$\times$
faster ops in a single PR); the rule matters for small
$<$\,5\,\% improvements where noise easily disguises the truth.

\section{Microbenchmarks vs End-to-End}
\label{sec:perf_methodology:micro_vs_e2e}

Two scales of measurement:

\subsection{Microbenchmarks}
\label{ssec:perf_methodology:micro}

Single-op timing: \code{lucid.matmul(a, b)} for a fixed shape,
or \code{F.softmax(x, dim=-1)}. Used for
\begin{itemize}
  \item Validating that a new emitter beats the composite path.
  \item Comparing dispatch overhead per call.
  \item Validating allocator behaviour under repeated
    same-shape allocations.
\end{itemize}

Microbenchmarks live in \code{lucid/bench/micro/}. The harness
sweeps over shape, dtype, and contiguity, producing a CSV that
\code{tools/perf-track} ingests for trending.

\subsection{End-to-end benchmarks}
\label{ssec:perf_methodology:e2e}

Full training/inference iteration of a model. Used for
\begin{itemize}
  \item Tracking the model-zoo performance over releases.
  \item Validating that microbenchmark wins translate to
    real workload wins.
  \item Catching regressions that microbenchmarks miss
    (memory pressure, allocator thrash).
\end{itemize}

End-to-end benchmarks live in \code{lucid/bench/e2e/}.

The micro-vs-end-to-end gap is informative: a microbench win
that does not translate to end-to-end usually means the op was
not on the critical path. We try to make both scales agree.

\section{The Performance Tracker}
\label{sec:perf_methodology:perf_tracker}

CI runs the full benchmark suite on every release-candidate
build. Results are committed to
\code{lucid/bench/results/<version>/`} as CSV. The
\code{tools/perf-track plot} command compares two releases:

\begin{lstlisting}[style=lucid,language=LucidPython]
$ tools/perf-track plot --baseline 3.4.0 --candidate 3.5.0-rc1
Regression on ssd_300_train: 130 ms -> 145 ms (+11.5%). FAIL.
Improvement on bert_base_train: 165 ms -> 140 ms (-15.2%).
Improvement on gpt2_small_train: 175 ms -> 95 ms (-45.7%).
\end{lstlisting}

Regressions $>5$\,\% block the release; the author of the
suspect PR investigates.

\section{Memory Profiling}
\label{sec:perf_methodology:memory_profiling}

Memory is the second metric (after time) that the bench harness
tracks. The \code{lucid.profiler.memory} module records:

\begin{itemize}
  \item Peak allocated bytes (the headline metric).
  \item Number of allocator calls (high count $\to$ allocator
    thrash).
  \item Allocator hit rate (cache hits vs miss-and-allocate).
  \item Per-op memory deltas (forward attribution).
\end{itemize}

\begin{lstlisting}[style=lucid,language=LucidPython]
with lucid.profiler.memory.profile() as mp:
    train_step(...)

print(mp.peak_bytes() / 1e9, "GB peak")
print(mp.allocator_calls(), "allocator calls")
\end{lstlisting}

\section{Energy Measurement}
\label{sec:perf_methodology:energy}

Power-aware benchmarks are reported in J/iteration when
relevant (mobile deployment, sustained training). The
\code{powermetrics} sampling at 100\,ms is averaged over the
benchmark window; CPU and GPU rails are aggregated.

For a typical M3~Max training run (ResNet-50, batch 32,
sustained), the power draw is roughly:
\begin{itemize}
  \item GPU: 32\,W average.
  \item CPU (used for data loading): 8\,W.
  \item Other (memory controller, SoC fabric): 4\,W.
  \item Total: 44\,W $\to$ at 78\,ms/iter that's 3.4\,J/iter.
\end{itemize}

The energy per training token / image is the metric for
deployment comparisons.

\section{Apple Silicon-Specific Notes}
\label{sec:perf_methodology:apple_specific}

\subsection{The chiplet boundary}
\label{ssec:perf_methodology:chiplet}

M-series chips have multiple GPU clusters; intra-cluster
operations are faster than inter-cluster. The MPSGraph dispatcher
handles cluster placement automatically, but particularly small
ops can be sensitive: a tensor that the dispatcher splits across
clusters incurs an extra synchronisation. The bench harness
notes when this is suspected (large variance + low utilisation).

\subsection{The performance-vs-efficiency cores}
\label{ssec:perf_methodology:p_vs_e_cores}

Apple silicon has both P-cores (high-performance) and E-cores
(low-power). Python's GIL bottleneck typically keeps the
\Lucid\ driver thread on one core; macOS's scheduler may pin it
to an E-core under low load. The bench harness sets the thread
QoS to \texttt{NSQualityOfServiceUserInteractive} to encourage
P-core placement.

\subsection{Memory pressure}
\label{ssec:perf_methodology:memory_pressure}

Above 80\% of unified memory utilisation, macOS begins paging.
A benchmark that crosses this threshold sees catastrophic
slowdowns. The bench harness tracks \code{vm\_stat} and warns
if pressure crosses the threshold during the run.

\section{Reproducibility of Results}
\label{sec:perf_methodology:reproducibility}

For every reported number, we ship enough information for
reproduction:

\begin{itemize}
  \item Git commit SHA of \Lucid\ at the time of measurement.
  \item Hardware configuration (sysctl output).
  \item macOS version.
  \item Power profile.
  \item Exact bench command.
\end{itemize}

The \code{lucid/bench/results/<version>/manifest.json} captures
all of these for each measurement set.

\section{Summary and Forward References}
\label{sec:perf_methodology:summary}

The methodology underpinning the performance work: synchronise
around timing, control for thermal drift, report median + p05/p95,
compare to reference-framework baseline on the same hardware. The
performance tracker enforces no-regression in CI; the energy
metric supports power-aware comparisons.

The next chapters apply the methodology:
\chref{ch:stabilization} on the stabilisation pass that got
\Lucid\ to within performance parity of the reference framework;
\chref{ch:amp_integration} on AMP performance; \chref{ch:mpsgraph_dispatch}
on MPSGraph dispatch tuning; \chref{ch:no_compile_sweep} on the
no-compile baseline; \chref{ch:graph_capture_wins} on what
compilation buys.

\begin{lucidnote}
For the user: when reporting an issue about performance, please
provide bench-harness output (the per-iter median, p05/p95, and
the reproduction command). ``Slower than \texttt{ref}'' alone is
not actionable; a 30-iter bench is.
\end{lucidnote}
